% Full title: Introduction
\section{Introduction}

Consider a dielectric structure that is periodic in one coordinate, extends
to infinity in the transverse coordinate, and differs from a perfect periodic
medium only in a finite-width central region.  At a fixed axial Bloch
quasimomentum, a guided mode is an eigenfunction that is square-integrable in
the transverse direction.  For the identical-background setting treated
there, the exact reduction of this infinite-domain eigenproblem to a bounded
defect region by Dirichlet-to-Neumann (DtN) operators is established in
\cite[Theorem~4.5]{Fliss2013}; the associated gap
eigenfunctions are exponentially localized by
\cite[Theorem~3.5]{Fliss2013}.  Propagation-operator and Riccati constructions
for periodic outlets were developed in \cite{JolyLiFliss2006,Coatleven2012},
and Robin-to-Robin maps provide a regular alternative at frequencies where a
Dirichlet chart fails \cite{FlissKlindworthSchmidt2015}.

The purpose of this draft is not to reprove that bounded-domain reduction.
Instead, it prepares a precise continuous framework in which the center cell is
represented by a M\"uller-type boundary integral ansatz while each half-guide is
represented by its complete Cauchy relation.  The latter is a linear relation
between Dirichlet and Neumann traces; it reduces to a DtN graph when its
Dirichlet projection is invertible.  The modal coordinates must include
generalized Floquet modes and Jordan chains.  Ordinary Bloch eigenvectors alone
need not be complete, whereas Riesz-basis and translation-operator results are
available in \cite[Theorem~2.2]{HohageSoussi2013}; related spectral
decompositions and modal DtN approximations are developed in
\cite{Zhang2021,Zhang2023}.

The center representation follows the quasiperiodic/free-space architecture
of \cite[Theorem~4 and Appendix~A]{BarnettGreengard2010}, but it is deliberately
stated on a quotient.  A physical field may be unique even when its layer
densities are not.  This distinction is essential because boundary integral
operators can exhibit representation-induced nullspaces or large nonphysical
inverse norms \cite{HiptmairMoiolaSpence2022}.  In particular, second-kind
structure alone does not imply that every small singular value represents a
guided mode.

The paper has four mathematical stages.  \Cref{sec:functional-setting}
defines the fixed-$\beta$ strip eigenproblem, records the two-ended
essential-spectrum formula as a proved theorem, and separates common
gaps from the compact gap windows used later.  \Cref{sec:cauchy-reduction}
defines weak Cauchy traces,
outgoing relations, their generalized-Bloch coordinates, and the exact
restriction--gluing reduction.  \Cref{sec:muller-coupling} fixes all
layer-potential conventions, introduces the center reconstruction and the
coupled operator, and isolates the representation nullspace.
\Cref{sec:main-results} states the quotient-safe main target and the remaining
Fredholm question.  Formal statements are followed by a \emph{proof roadmap}
only when a proof or adaptation remains to be completed.
